\documentclass[12pt, a4paper]{article} % Classe: article, report o book [8]

% --- Pacchetti Scientifici/Matematici ---
\usepackage{amsmath, amssymb, amsfonts, amsthm} % Matematica avanzata [2]
\usepackage{booktabs} % Per tabelle professionali

% --- Formattazione Testo e Altro ---
\usepackage[document]{ragged2e}
\usepackage{xcolor}    % Per usare i colori 

% --- Codifica e Font ---
\usepackage[utf8]{inputenc}
\usepackage[T1]{fontenc}
\usepackage[italian]{babel}

% Scelta del font (Times)
\usepackage{mathptmx} 

% --- Gestione Margini e Layout (Standard 65-70 caratteri) ---
\usepackage[includeheadfoot]{geometry}
\geometry{
    a4paper,
    top=3cm,
    bottom=3cm,
    left=3.5cm,   % Margine più ampio per la rilegatura
    right=3.5cm,  % Regolando questi margini si ottengono ~65-70 caratteri
}

% --- Gestione Righe (32-35 righe per pagina) ---
\usepackage{setspace}
\setstretch{1.3} % Interlinea personalizzata per circa 33-34 righe

% --- Gestione Figure e Tavole ---
\usepackage{graphicx}
\usepackage[figuresright]{rotating} % Per tavole in A3 o orizzontali

% --- PDF e Hyperlinks ---
\usepackage[pdftex, hidelinks]{hyperref}

% Document metadata
\title{Risultato delle interviste}
\author{Mattia Cavina}
\date{03 febbraio 2026}

\begin{document}
Il presente progetto di tesi riguarda la progettazione e implementazione di un nuovo configuratore commerciale per l'azienda Cepi~\cite{Cepi}. 
Attualmente l'azienda utilizza un programma sviluppato internamente basato su Access, il quale presenta significative limitazioni tecniche e funzionali dopo circa 30 anni di utilizzo. 
L'analisi si concentra sulla valutazione del sistema attuale, sulla raccolta sistematica dei requisiti attraverso interviste con gli utenti, e sulla progettazione di un nuovo database relazionale.
Il progetto è stato realizzato in collaborazione con il project manager esterno Claudio Rossi, che ha coordinato l'avanzamento dei lavori e ha contribuito all'implementazione della soluzione.
\end{document}